% !TeX root = ../main.tex

% ==================================================
% A.13. 실전 응용·예외·통합 정리
% ==================================================

\subsection{실전 응용·예외·통합 정리}

Function Object · Name Binding · Closure · Decorator · Descriptor · Registration · Exception Handling · Async

\subsubsection{이번 학습의 목표}

\zsh{A.1}부터 \zsh{A.12}까지 Python의 함수 객체와 데코레이터를 여러 관점에서 분석하였다.

처음에는 함수 저으이와 호출의 차이에서 출발했다.

\begin{dhkcode}
def calculate(x):
    return x * 2
\end{dhkcode}

이후 중첩 함수와 클로저를 학습했고, 함수 객체를 다른 함수에 전달하여 새로운 함수를 반환하는 데코레이터를 다뤘다.

\begin{dhkcode}
@trace
def calculate(x):
    return x * 2
\end{dhkcode}

\zsh{A.11}에서는 클래스 네임스페이스에 저장된 함수 객체가 디스크립터 프로토콜을 통해 바운드 메서드로 변환되는 과정을 분석했다.

마지막으로 \zsh{A.12}에서는 FastAPI의 등록형 데코레이터를 통해 함수 객체가 외부 시스템에 등록되고, 나중에 HTTP 요청을 통해 호출되는 구조까지 살펴보았다.

이제 남은 과제는 이 개념들을 하나의 일관된 Python 실행 모델로 통합하는 것이다.

이번 학습에서는 새로운 문법을 많이 추가하지 않을 것이다.

대신 지금까지 배운 개념이 실제 프로그램에서 어떻게 상호작용하는지를 추적할 것이다.

핵심 질문은 다음과 같다.

\begin{dhksmry}[centered]
여러 데코레이터가 적용된 클래스 메서드를 호출할 때,

어떤 객체가 생성되고, 어떤 이름이 어디에 바인딩되며,

어떤 디스크립터와 클로저를 거쳐 최종 함수가 실행될까?
\end{dhksmry}

그리고 정상적인 호출뿐 아니라 예외 발생, 상태 변경, 비동기 함수, 등록형 데코레이터의 경계 조건까지 분석할 것이다.

\subsubsection{통합 분석을 위한 예제}

이번에는 다음 기능을 가진 작은 시스템을 구현해 보자.

\begin{dhksmry}
- register: 함수 객체를 외부 레지스트리에 등록하는 데코레이터 팩토리

- monitor: 함수 호출 횟수를 기록하는 데코레이터 팩토리

- classmethod: 함수를 클래스 객체에 바인딩하는 디스크립터

- Base, Derived: 상속과 MRO를 확인하기 위한 클래스
\end{dhksmry}

이 예제 하나로 \zsh{A.1}부터 \zsh{A.12}까지의 핵심 개념 대부분을 연결할 수 있다.

\begin{dhknote}
\begin{dhksmry}[centered]
전체 코드
\end{dhksmry}
\end{dhknote}

\begin{dhkcode}
from functools import wraps


registry = {}


def register(key):
    print("evaluate register", key)

    def decorator(func):
        print("apply register", key)

        registry[key] = func

        return func

    return decorator


def monitor(tag):
    print("evaluate monitor", tag)

    def decorator(func):
        print("apply monitor", tag)

        calls = 0

        @wraps(func)
        def wrapper(*args, **kwargs):
            nonlocal calls
            calls += 1

            print("enter", tag, calls)

            try:
                return func(*args, **kwargs)
            finally:
                print("leave", tag)

        wrapper.count = lambda: calls

        return wrapper

    return decorator


class Base:
    scale = 10

    @classmethod
    @register("predict")
    @monitor("predict")
    def predict(cls, value):
        print("body", cls.__name__, value)
        return cls.scale * value


class Derived(Base):
    scale = 20


print("ready")

print(
    "identity",
    registry["predict"]
    is Base.__dict__["predict"].__func__
)

print(
    "binding",
    Derived.predict.__self__ is Derived
)

print("result", Derived.predict(3))

try:
    Derived.predict(None)
except TypeError:
    print("caught TypeError")

print("calls", registry["predict"].count())
\end{dhkcode}

\begin{dhknote}
\begin{dhksmry}[centered]
실행 결과
\end{dhksmry}
\end{dhknote}

다음 출력은 Python에서 실행하여 확인할 수 있는 결과다.

\begin{dhkzsh}
evaluate register predict
evaluate monitor predict
apply monitor predict
apply register predict
ready
identity True
binding True
enter predict 1
body Derived 3
leave predict
result 60
enter predict 2
body Derived None
leave predict
caught TypeError
calls 2
\end{dhkzsh}

여기서 우리가 설명해야 할 현상은 여섯 가지다.

\begin{dhksmry}
왜 evaluate register가 evaluate monitor보다 먼저 출력될까?

왜 적용 순서는 반대로 apply monitor가 먼저일까?

왜 레지스트리에 저장된 함수와 classmethod.__func__가 동일할까?

왜 Base에 정의된 함수의 cls가 Derived를 참조할까?

왜 Derived.predict(None)에서 예외가 발생해도 leave가 출력될까?

왜 정상 호출은 한 번인데 호출 횟수는 2일까?
\end{dhksmry}

\subsubsection{첫 번째 단계 - 데코레이터 표현식 평가}

가장 먼저 다음 부분을 살펴보자.

\begin{dhkcode}
class Base:
    scale = 10

    @classmethod
    @register("predict")
    @monitor("predict")
    def predict(cls, value):
        ...
\end{dhkcode}

Python은 클래스 본문을 실행하는 과정에서 \zsh{predict} 함수 정의문을 처리한다.

이때 데코레이터 표현식은 위에서 아래로 평가된다.

따라서 다음 순서가 성립한다.

\begin{dhksmry}
    1. classmethod 평가

    2. register("predict") 호출

    3. monitor("predict") 호출

    4. 원본 predict 함수 객체 생성
\end{dhksmry}

\zsh{classmethod}는 이미 존재하는 내장 객체를 참조하므로 별도의 출력이 없다.

반면 \zsh{register("predict")}는 팩토리 함수를 호출하므로 다음 출력이 발생한다.

\begin{dhkzsh}
evaluate register predict
\end{dhkzsh}

그리고 \zsh{monitor("predict")}를 호출하면서 다음 출력이 발생한다.

\begin{dhkzsh}
evaluate monitor predict
\end{dhkzsh}

여기서 중요한 사실이 있다.

아직 원본 \zsh{predict} 함수 본문은 실행되지 않았다.

\zsh{register}와 \zsh{monitor}는 데코레이터 팩토리로서 실행되었지만, \zsh{predict}는 호출된 것이 아니라 함수 객체로 생성되는 과정에 있다.

따라서 \zsh{body}는 출력되지 않는다.

이전 단원과 연결하면 다음과 같다.

\begin{dhksmry}
    A.1: 함수 객체와 함수 호출은 다르다.

    A.2: 중첩 함수는 함수 객체를 반환할 수 있다.

    A.9: 데코레이터 팩토리는 데코레이터 함수 객체를 반환한다.

    A.10: 여러 데코레이터 표현식은 위에서 아래로 평가된다.
\end{dhksmry}

\subsubsection{두 번째 단계 - 데코레이터 적용과 클래스 네임스페이스}

데코레이터 표현식의 평가가 끝나면 원본 \zsh{predict} 함수 객체가 생성되고, 데코레이터가 아래에서 위로 적용된다.

개념적으로 다음과 대응한다.

\begin{dhkcode}
register_decorator = register("predict")
monitor_decorator = monitor("predict")

original = predict

wrapper = monitor_decorator(original)

registered = register_decorator(wrapper)

descriptor = classmethod(registered)

predict = descriptor
\end{dhkcode}

이 코드는 객체 관계를 이해하기 위한 개념적 변환식이다.

실제 클래스 정의문의 내부 실행 절차를 문장 단위로 완전히 재현한 것은 아니다.

\begin{dhknote}
\begin{dhksmry}[centered]
1. monitor 적용
\end{dhksmry}
\end{dhknote}

먼저 다음 호출이 수행된다.

\begin{dhkcode}
wrapper = monitor_decorator(original)
\end{dhkcode}

\zsh{monitor_decorator}는 \zsh{calls = 0}을 초기화하고 새로운 wrapper 함수 객체를 생성한다.

이 wrapper는 클로저를 토앻 다음 객체를 참조할 수 있다.

\begin{dhkzsh}
calls → 정수 객체 0
func  → 원본 predict 함수 객체
tag   → 문자열 "predict"
\end{dhkzsh}

그리고 wrapper 함수 객체를 반환한다.

이 과정에서 다음 출력이 발생한다.

\begin{dhkzsh}
apply monitor predict
\end{dhkzsh}

\begin{dhknote}
\begin{dhksmry}[centered]
2. register ㅈ거용
\end{dhksmry}
\end{dhknote}

다음 호출이 수행된다.

\begin{dhkcode}
registered = register_decorator(wrapper)
\end{dhkcode}

이때 \zsh{register}에 전달되는 것은 원본 함수가 아니다.

이미 \zsh{monitor}가 반환한 wrapper 함수 객체다.

따라서 다음 문장은

\begin{dhkcode}
registry[key] = func
\end{dhkcode}

개념적으로 다음과 같다.

\begin{dhkcode}
registry["predict"] = wrapper
\end{dhkcode}

그리고 \zsh{register}는 동일한 wrapper 객체를 반환한다.

따라서 다음 출력이 발생한다.

\begin{dhkzsh}
apply register predict
\end{dhkzsh}

\begin{dhknote}
\begin{dhksmry}[centered]
3. classmethod 적용
\end{dhksmry}
\end{dhknote}

마지막으로 다음 호출이 수행된다.

\begin{dhkcode}
descriptor = classmethod(wrapper)
\end{dhkcode}

이 호출은 wrapper 함수 객체를 감싸는 \zsh{classmethod} 디스크립터를 생성한다.

따라서 \zsh{Base} 클래스 네임스페이스에는 원본 함수도, wrapper 함수도 직접 저장되지 않는다.

최종적으로 \zsh{classmethod} 객체가 저장된다.

전체 구조는 다음과 같다.

\begin{dhkgraph}[uniform]
node n1 (1,1) :: Base.__dict__["predict"] // classmethod 디스크립터
node n2 (2,1) :: monitor의 wrapper 함수 객체 // 클로저: calls, func, tag // registry["predict"]도 이 객체 참조
node n3 (3,1) :: 원본 predict 함수 객체 // cls.scale * value 실행

edge n1 -> n2 :: __func__
edge n2 -> n3 :: func / __wrapped__

explain :: 등록·wrapper·디스크립터는 서로 다른 객체 관계
\end{dhkgraph}

이 구조를 정확히 이해하면 \zsh{A.11}과 \zsh{A.12}가 하나로 연결된다.

\begin{dhksmry}
    - monitor는 wrapper를 생성한다.

    - register는 wrapper를 외부 레지스트리에 등록한다.

    - classmethod는 wrapper의 클래스 바인딩을 담당한다.
\end{dhksmry}

세 데코레이터의 역할은 서로 다르다.

\subsubsection{세 번째 단계 - 속성 탐색과 클래스 메서드 바인딩}

클래스 정의가 완료되었으므로 이제 다음 표현식을 평가할 수 있다.

\begin{dhkcode}
Derived.predict
\end{dhkcode}

\zsh{Derived}는 \zsh{predict}를 직접 정의하지 않았다.

따라서 클래스 속성 탐색은 MRO를 따라 \zsh{Base}에서 \zsh{predict}를 찾는다.

\begin{dhkcode}
print(Derived.__mro__)
\end{dhkcode}

예상 출력:

\begin{dhkzsh}
(<class '__main__.Derived'>, <class '__main__.Base'>, <class 'object'>)
\end{dhkzsh}

\zsh{Base}에서 발견되는 객체는 \zsh{classmethod} 디스크립터다.

이 디스크립터는 자신이 감싸고 있는 wrapper 함수 객체를 접근한 클래스인 \zsh{Derived}에 바인딩한다.

따라서 다음과 같은 바운드 메서드가 만들어진다.

\begin{dhkdir}[in=6]
Derived.predict
|바운드 메서드
> __self__ -> Derived 클래스 객체
> __func__ -> monitor의 wrapper 함수 객체
\end{dhkdir}

이 관계는 다음 코드로 검증할 수 있다.

\begin{dhkcode}
method = Derived.predict

print(method.__self__ is Derived)

print(
    method.__func__
    is registry["predict"]
)
\end{dhkcode}

예상 출력:

\begin{dhkzsh}
True
True
\end{dhkzsh}

여기서 중요한 점은 두 가지다.

첫째, \zsh{predict}를 정의한 클래스와 실제 바인딩되는 클래스는 다를 수 있다.

\zsh{predict}는 \zsh{Base}에서 정의되었지만, 

\zsh{Derived.predict}를 조회했으므로 \zsh{__self__} 는 \zsh{Derived} 를 참조하다.

둘째, 바인딩되는 함수는 원본 \zsh{predict}가 아니라 \zsh{monitor}가 반환한 wrapper다.

\zsh{classmethod}는 자신이 전달받은 ㅈrapper를 클래스에 바인딩할 뿐, 그 wrapper가 클로저로 참조하는 원본 함수까지 별도로 바인딩하지는 않는다.

원본 함수는 나중에 wrapper 내부에서 호출된다.

\subsubsection{네 번째 단계 - 실제 함수 호출과 반환}

이제 다음 표현식을 실행해 보자.

\begin{dhkcode}
Derived.predict(3)
\end{dhkcode}

이 호출은 개념적으로 다음과 같다.

\begin{dhkcode}
wrapper(Derived, 3)
\end{dhkcode}

바운드 메서드가 \zsh{Derived} 클래스 객체를 첫 번째 인수로 전달하기 때문이다.

\begin{dhknote}
\begin{dhksmry}[centered]
1. Wrapper 호출 프레임 생성
\end{dhksmry}
\end{dhknote}

Wrapper의 매개변수는 다음과 같다.

\begin{dhkcode}
def wrapper(*args, **kwargs):
    ...
\end{dhkcode}

따라서 다음 바인딩이 형성된다.

\begin{dhkzsh}
args = (Derived, 3)
kwargs = {}
\end{dhkzsh}

여기서 \zsh{args}는 튜플 객체이고, \zsh{kwargs}는 딕셔너리 객체다.

\zsh{Derived} 클래스 객체나 정수 3이 복사되어 새로운 객체가 만들어진다는 뜻은 아니다.

이 객체들에 대한 참조가 인수 전달 구조에 포함되는 것이다.

\begin{dhknote}
\begin{dhksmry}[centered]
2. 클로저 상태 변경
\end{dhksmry}
\end{dhknote}

다음 코드가 실행된다.

\begin{dhkcode}
nonlocal calls
calls += 1
\end{dhkcode}

\zsh{calls}는 wrapper 호출 프레임의 지역 변수가 아니다.

바깥쪽 \zsh{decorator} 함수의 실행에서 생성된 클로저 셀에 연결된 이름이다.

처음에는 \zsh{calls}가 정수 객체 0을 참조했다.

\zsh{calls += 1}이 실행되면 새로운 정수 결과 객체에 \zsh{calls}가 재바인딩된다.

\begin{dhksmry}
호출 전

calls 셀 \ra 0

호출 후

calls 셀 \ra 1
\end{dhksmry}

정수 객체 자체를 변경하는 것이 아니다.

이후 다음 문장이 실행된다.

\begin{dhkcode}
print("enter", tag, calls)
\end{dhkcode}

따라서 다음 출력이 발생한다.

\begin{dhkzsh}
enter predict 1
\end{dhkzsh}

\begin{dhknote}
\begin{dhksmry}[centered]
3. 원본 함수 호출
\end{dhksmry}
\end{dhknote}

다음 문장이 실행된다.

\begin{dhkcode}
return func(*args, **kwargs)
\end{dhkcode}

\zsh{func}은 클로저를 통해 원본 \zsh{predict} 함수 객체를 참조한다.

따라서 다음과 같은 호출이 수행된다.

\begin{dhkcode}
original_predict(Derived, 3)
\end{dhkcode}

원본 함수의 매개변수 바인딩은 다음과 같다.

\begin{dhkzsh}
cls = Derived
value = 3
\end{dhkzsh}

이제 원본 함수 본문이 실행된다.

\begin{dhkcode}
print("body", cls.__name__, value)

return cls.scale * value
\end{dhkcode}

먼저 다음 출력이 발생한다.

\begin{dhkzsh}
body Derived 3
\end{dhkzsh}

그리고 다음 표현식이 평가된다.

\begin{dhkcode}
cls.scale
\end{dhkcode}

\zsh{cls}는 \zsh{Derived} 클래스 객체를 참조한다.

\zsh{Dervied}의 클래스 네임스페이스에는 \zsh{scale = 20}이 정의되어 있으므로 그 값을 사용한다.

따라서 계산 결과는 다음과 같다.

\[
20 \times 3 = 60
\]

원본 함수는 정수 객체 60을 반환한다.

\begin{dhknote}
\begin{dhksmry}[centered]
4. finally 실행과 반환
\end{dhksmry}
\end{dhknote}

Wrapper의 코드를 다시 보자.

\begin{dhkcode}
try:
    return func(*args, **kwargs)
finally:
    print("leave", tag)
\end{dhkcode}

\zsh{return} 문장이 실행되더라도 함수가 즉시 종료되는 것은 아니다.

\zsh{try}블록에서 반환값이 결정된 뒤, 해당 \zsh{finally} 블록이 실행된다.

따라서 다음 출력이 발생한다.

\begin{dhkzsh}
leave predict
\end{dhkzsh}

이후 반환값 60이 호출자에게 전달된다.

마지막으로 바깥쪽 \zsh{print}가 실행되면서 다음 출력이 발생한다.

\begin{dhkzsh}
result 60
\end{dhkzsh}

\begin{dhknote}
\begin{dhksmry}[centered]
전체 호출 흐름
\end{dhksmry}
\end{dhknote}

\begin{dhkgraph}[uniform]
node n1 (1,1) :: Derived.predict(3) // 바운드 메서드 호출
node n2 (2,1) :: wrapper(Derived, 3) // calls:0 -> 1
node n3 (3,1) :: original_predict(Derived, 3) // cls.scale = 20
node n4 (4,1) :: 60 반환 // finally 실행: leave predict
node n5 (5,1) :: 최종 결과: 60

edge n1 -> n2
edge n2 -> n3
edge n3 -> n4

explain :: Derived.predict(3)의 호출 및 반환 경로

caption :: 메서드 바인딩은 wrapper 호출 이전에, 클로저 탐색은 wrapper 실행 중에 발생한다.
\end{dhkgraph}

이 과정에서 우리가 학습한 객체 모델이 모두 연결된다.

\begin{dhksmry}
    __get__: 클래스 바인딩

    __self__: 바인딩된 Derived

    __func__: wrapper 함수 객체

    *args, **kwargs: 인수 전달

    __closure__: 원본 함수와 상태 참조

    nonlocal: 클로저 셀의 이름 재바인딩

    return: 원본 함수의 반환값 전달

    finally: 반환 전 정리 코드 실행
\end{dhksmry}

\subsubsection{다섯 번째 단계 - 예외가 발생하면 무엇이 달라질까?}

이제 다음 코드를 실행한다고 하자.

\begin{dhkcode}
try:
    Derived.predict(None)
except TypeError:
    print("caught TypeError")
\end{dhkcode}

앞의 정상 호출과 같은 경로를 따라 실행되지만, 원본 함수의 연산에서 예외가 발생한다.

\begin{dhknote}
\begin{dhksmry}[centered]
1. Wrapper 실행
\end{dhksmry}
\end{dhknote}

\zsh{erived.predict(None)}을 조회하고 호출하면 wrapper는 다음 인수를 전달받는다.

\begin{dhkzsh}
args = (Derived, None)
\end{dhkzsh}

이후 \zsh{calls += 1}이 실행된다.

첫 번째 호출에서 \zsh{calls}가 1이 되었으므로 이번에는 2가 된다.

\begin{dhkzsh}
enter predict 2
\end{dhkzsh}

\begin{dhknote}
\begin{dhksmry}[centered]
2. 원본 함수 실행
\end{dhksmry}
\end{dhknote}

원본 함수는 다음 바인딩으로 실행된다.

\begin{dhkzsh}
cls = Derived
value = None
\end{dhkzsh}

따라서 다음 출력이 발생한다.

\begin{dhkzsh}
body Derived None
\end{dhkzsh}

이후 다음 표현식이 평가된다.

\begin{dhkcode}
cls.scale * value
\end{dhkcode}

이는 다음과 같다.

\begin{dhkcode}
20 * None
\end{dhkcode}

정수와 \zsh{None} 사이에는 이 곱셈을 수행하는 연산이 정의되어 있지 않으므로 \zsh{TypeError}가 발생한다.

\begin{dhknote}
\begin{dhksmry}[centered]
3. 예외 전파와 finally
\end{dhksmry}
\end{dhknote}

원본 함수가 정상적으로 반환하지 못했으므로 wrapper의 \zsh{return}도 정상적으로 완료되지 못한다.

하지만 \zsh{finally} 블록은 실행된다.

\begin{dhkzsh}
leave predict
\end{dhkzsh}

Wrapper는 예외를 처리하는 \zsh{except} 블록을 갖고 있지 않으므로 \zsh{TypeError}는 호출자에게 전파된다.

바깥쪽 \zsh{except TypeError}가 이를 처리하고 다음 내용을 출력한다.

\begin{dhkzsh}
caught TypeError
\end{dhkzsh}

\begin{dhknote}
\begin{dhksmry}[centered]
4. 클로저의 상태는 어떻게 될까?
\end{dhksmry}
\end{dhknote}

여기서 중요한 점이 있다.

함수 호출은 실패했지만 다음 코드는 이미 실행되었다.

\begin{dhkcode}
calls += 1
\end{dhkcode}

따라서 클로저 상태는 이전 값으로 자동 복구되지 않는다.

\begin{dhkcode}
print(registry["predict"].count())
\end{dhkcode}

결과는 다음과 같다.

\begin{dhkzsh}
2
\end{dhkzsh}

이것은 일반적인 Python 예외 처리의 중요한 성질이다.

\begin{dhksmry}[centered]
예외가 발생했다고 해서 예외 이전에 수행된 모든 상태 변경이

자동으로 취소되는 것은 아니다.
\end{dhksmry}

이번 예제의 \zsh{calls}는 성공 횟수가 아니라 wrapper 진입 횟수를 나타낸다.

성공한 호출만 세고 싶다면 상태를 변경하는 위치를 바꿔야 한다.

예를 들어 다음과 같이 설계할 수 있다.

\begin{dhkcode}
result = func(*args, **kwargs)
calls += 1
return result
\end{dhkcode}

이 경우 원본 함수가 정상적으로 반환한 뒤에만 \zsh{calls}를 증가시킨다.

다만 wrapper에서 이후 수행하는 처리까지 모두 성공해야 성공으로 볼 것인지에 따라 더 세밀한 설계가 필요할 수 있다.

\subsubsection{클로저 셀과 함수 메타데이터를 직접 검증하기}

이제 앞의 통합 예제를 객체 내부 구조 관점에서 검증해 보자.

먼저 클래스 네임스페이스에 저장된 \zsh{classmethod} 객체를 가져온다.

\begin{dhkcode}
descriptor = Base.__dict__["predict"]
\end{dhkcode}

그 객체가 감싸고 있는 wrapper를 가져온다.

\begin{dhkcode}
wrapper = descriptor.__func__
\end{dhkcode}

그리고 \zsh{functools.wraps}가 설정한 \zsh{__wrapped__}를 이용하여 원본 함수 객체를 확인한다.

\begin{dhkcode}
original = wrapper.__wrapped__
\end{dhkcode}

다음 코드를 실행해 보자.

\begin{dhkcode}
print(type(descriptor).__name__)
print(type(wrapper).__name__)

print(wrapper is registry["predict"])

print(
    original is wrapper.__wrapped__
)
\end{dhkcode}

예상 출력:

\begin{dhkzsh}
classmethod
function
True
True
\end{dhkzsh}

\begin{dhknote}
\begin{dhksmry}[centered]
__closure__와 __wrapped__는 다르다
\end{dhksmry}
\end{dhknote}

\zsh{__wrapped__}는 \zsh{functools.wraps}가 설정한 함수 메타데이터 참조다.

반면 \zsh{__closure__}는 함수 객체가 참조하는 클로저 셀들의 튜플이다.

두 구조는 서로 다른 목적을 가진다.

\begin{dhkcode}
def closure_map(func):
    names = func.__code__.co_freevars
    cells = func.__closure__ or ()

    return dict(zip(names, cells))


cells = closure_map(wrapper)

print(wrapper.__code__.co_freevars)

print(cells["calls"].cell_contents)

print(
    cells["func"].cell_contents
    is original
)
\end{dhkcode}

앞의 통합 예제를 모두 실행한 뒤라면 결과는 다음과 같다.

\begin{dhkzsh}
('calls', 'func', 'tag')
2
True
\end{dhkzsh}

여기서 \zsh{co_freevars}는 함수 코드 객체가 사용하는 자유 변수 이름들을 보여준다.

\zsh{__closure__}는 그 이름에 대응하는 셀 객체들을 제공한다.

따라서 \zsh{closure_map}을 이용하면 이름과 셀을 연결해서 확인할 수 있다.

이 예제에서는 다음 관계가 형성된다.

\begin{dhkdir}[in=6]
wrapper 함수 객체
> __code__
>> co_freevars
>>> ("calls", "func", "tag")
> __closure__
>> calls 셀 -> 2
>> func 셀 -> 원본 predict 함수
>> tag 셀 -> "predict"
> __wrapped__ -> 원본 predict 함수
\end{dhkdir}

\zsh{func} 셀과 \zsh{__wrapped__}는 이 예제에서 같은 원본 함수 객체를 참조한다.

하지만 둘이 같은 개념인 것은 아니다.

\zsh{func} 셀은 wrapper가 실행될 때 원본 함수를 호출하기 위해 사용하고,

\zsh{__wrapped__}는 도구가 원본 함수를 추적할 수 있도록 제공되는 메타데이터다.

\begin{dhknote}
\begin{dhksmry}[centered]
wrapper.count는 어떻게 같은 상태를 읽을까?
\end{dhksmry}
\end{dhknote}

다음 문장을 기억해 보자.

\begin{dhkcode}
wrapper.count = lambda: calls
\end{dhkcode}

이 코드는 wrapper 함수 객체의 \zsh{__dict__}에 \zsh{"count"} 라는 속성을 추가한다.

그 속성값은 \zsh{lambda}로 생성한 별도의 함수 객체다.

그런데 이 함수 역시 같은 \zsh{calls} 셀을 참조한다.

다음과 같이 확인할 수 있다.

\begin{dhkcode}
wrapper_cells = closure_map(wrapper)
count_cells = closure_map(wrapper.count)

print(
    wrapper_cells["calls"]
    is count_cells["calls"]
)
\end{dhkcode}

예상 출력:

\begin{dhkzsh}
True
\end{dhkzsh}

즉, 두 함수 객체는 서로 다른 코드와 호출 동작을 갖지만 같은 클로저 셀을 공유한다.

\begin{dhkdiagram}[in=6]
node n1 (1,1) :: wrapper 함수
node n2 (3,1) :: count 함수
node n3 (2,2) :: calls 셀 -> 2

edge n1 -> n3 :: HV
edge n2 -> n3 :: HV
\end{dhkdiagram}

따라서 wrapper에서 \zsh{nonlocal calls}로 상태를 변경하면

\zsh{wrapper.count()}가 그 변경된 값을 읽을 수 있다.

여기서도 핵심은 함수 객체의 동일성이 아니라 공유하는 클로저 셀의 동일성이다.

\subsubsection{\zsh{inspect.signature()}는 무엇을 보여줄까?}

\zsh{A.8}과 \zsh{A.11}의 내용을 연결해 보자.

앞의 통합 예제에 다음 코드를 추가한다.

\begin{dhkcode}
import inspect


print(inspect.signature(wrapper))

print(
    inspect.signature(
        wrapper,
        follow_wrapped=False
    )
)

print(inspect.signature(Derived.predict))
\end{dhkcode}

예상 출력:

\begin{dhkzsh}
(cls, value)
(*args, **kwargs)
(value)
\end{dhkzsh}

세 결과가 다른 이유를 분석해 보자.

\begin{dhktable}
|표현식|	결과|	원인|
|:---|:---|:---|
|inspect.signature(//*wrapper)|	(cls, value)|	__wrapped__를 따라 원본 함수의 시그니처 분석|
|inspect.signature(//*wrapper,//*follow_wrapped=False)|	(*args, **kwargs)|	실제 wrapper의 매개변수 구조 확인|
|inspect.signature(//*Derived.predict)|	(value)	|원본 시그니처와 클래스 바인딩 상태를 함께 고려|
\end{dhktable}

여기서 중요한 구별이 있다.

\zsh{@wraps}가 실제 wrapper의 매개변수를 \zsh{(cls, value)} 로 변경한 것은 아니다.

실제 wrapper의 매개변수는 여전히 \zsh{*args}, \zsh{**kwargs}다.

\zsh{inspect.signature}가 원본 함수의 메타데이터르 ㄹ추적하여 외부에 보여주는 시그니처를 구성하는 것이다.

그리고 \zsh{Derived.predict}는 클래스에 바인딩된 메서드이므로 호출자가 직접 전달할 필요가 없는 \zsh{cls}가 표시되지 않는다.

따라서 실제 함수 호출 규칙, 함수 메타데이터, 바운드 메서드의 시그니처 표현은 구별해야 한다.

\subsubsection{레지스트리에 저장된 함수는 자동으로 클래스에 바인딩될까?}

앞의 통합 예제에서 다음 관계가 성립했다.

\begin{dhkcode}
registry["predict"] is Base.__dict__["predict"].__func__
\end{dhkcode}

두 표현식은 동일한 wrapper 함수 객체를 참조한다.

그렇다면 다음 두 호출도 동일할까?

\begin{dhkcode}
Derived.predict(3)

registry["predict"](3)
\end{dhkcode}

그렇지 않다.

첫 번째는 \zsh{classmethod} 디스크립터를 통해 생성된 바운드 메서드의 호출이다.

따라서 \zsh{Derived} 클래스 객체가 자동으로 첫 번재 인수에 전달된다.

반면 두 번째는 딕셔너리에서 꺼낸 일반 함수 객체를 직접 호출하는 것이다.

딕셔너리 조회는 \zsh{classmethod} 바인딩을 수행하지 않는다.

따라서 wrapper에 전달되는 인수는 3 하나뿐이다.

Wrapper가 원본 함수에 그 인수를 전달하면

원본 \zsh{predict(cls, value)}에서 \zsh{value}에 

해당하는 인수가 부족하여 \zsh{TypeError} 가 발생한다.

레지스트리를 통해 정상적으로 호출하려면 다음처럼 클래스를 명시적으로 전달해야 한다.

\begin{dhkcode}
registry["predict"](Derived, 3)
\end{dhkcode}

이 호출에서는 wrapper가 다음 인수를 전달 받는다.

\begin{dhkzsh}
args = (Derived, 3)
\end{dhkzsh}

따라서 원본 함수도 정상적으로 실행할 수 있다.

이 차이는 중요한 판단 기준이다.

같은 함수 객체를 참조하더라도 어떤 경로로 접근했는지에 따라 디스크립터 바인딩이 적용될 수도 있고 적용되지 않을 수도 있다.

\zsh{A.11}에서 학습한 속성 탐색과 \zsh{A.12}에서 학습한 등록형 데코레이터가 만나는 지점이다.

\subsubsection{예외 처리에서 주의해야 할 \zsh{return}과 \zsh{finally}}

앞의 \zsh{monitor}는 다음과 같이 구현했다.

\begin{dhkcode}
try:
    return func(*args, **kwargs)
finally:
    print("leave", tag)
\end{dhkcode}

이 구조는 원본 함수가 정상 반환하든 예외를 발생시키든 정리 코드를 실행한다.

하지만 \zsh{finally} 블록에서 직접 \zsh{return}을 수행하면 의미가 달라진다.

\begin{dhkcode}
from functools import wraps


def unsafe(func):

    @wraps(func)
    def wrapper(*args, **kwargs):

        try:
            return func(*args, **kwargs)
        finally:
            return "forced"

    return wrapper


@unsafe
def fail():
    raise ValueError("failed")


print(fail())
\end{dhkcode}

예상 출력:

\begin{dhkzsh}
forced
\end{dhkzsh}

원본 함수에서 \zsh{ValueError}가 발생했지만 호출자에게 전달되지 않는다.

왜 그럴까?

\zsh{finally} 블록의 \zsh{return}이 이전에 진행 중이던 예외 전파를 중단하고 새로운 반환값을 결정했기 때문이다.

따라서 wrapper를 구현할 때는 다음 세 가지를 구별해야 한다.

\begin{dhktable}
|구현	|정상 반환|	예외 발생|
|:---|:---|:---|
|return func(...)|	원본 반환값 전달	|예외 전파|
|try: return func(...) + finally에서 로깅|	정리 코드 실행 후 반환	|정리 코드 실행 후 예외 전파|
|finally에서 직접 return|	finally의 값으로 반환	|예외가 억제될 수 있음|
\end{dhktable}

여기서 두 번재 경우에도 \zsh{finally} 내부에서 새로운 예외가 발생하면 기존 예외를 대체하는 등 결과가 달라질 수 있다.

따라서 로깅이나 리소스 정리를 위한 데코레이터에서는 예외 전파 규칙을 명확하게 설계해야 한다.

\subsubsection{비동기 함수와 제네레이터는 호출 시점이 다르다}

지금까지의 통합 예제는 일반적인 동기 함수를 대상으로 했다.

하지만 데코레이터를 비동기 함수나 제네레이터 함수에 적용할 때는 호출 결과와 실행 시점을 추가로 고려해야 한다.

\begin{dhknote}
\begin{dhksmry}[centered]
비동기 함수
\end{dhksmry}
\end{dhknote}

다음 코드를 보자.

\begin{dhkcode}
from functools import wraps
import inspect


def trace(func):

    @wraps(func)
    def wrapper(*args, **kwargs):
        print("wrapper")
        return func(*args, **kwargs)

    return wrapper


@trace
async def calculate(x):
    print("body")
    return x * 2


print(inspect.iscoroutinefunction(calculate))
\end{dhkcode}

예상 출력:

\begin{dhkzsh}
False
\end{dhkzsh}

원본 \zsh{calculate}는 코루틴 함수지만, 데코레이터가 반환한 객체는 일반적인 동기 wrapper 함수다.

\zsh{@wraps}는 wrapper를 코루틴 함수로 변환하지 않는다.

따라서 \zsh{calculate(5)}를 호출하면 동기 wrapper가 실행되고, 원본 코루틴 함수의 호출 결과인 코루틴 객체가 반환된다.

코루틴 본문은 그 객체가 실제로 실행될 때 진행된다.

FastAPI같은 프레임워크에서는 최종 등록된 함수가 동기 함수인지 비동기 함수인지가 실행 방식에 영향을 주기 때문에 중요하다.

일반적인 코루틴 함수와 동기 함수를 모두 지원하려면 다음처럼 wrapper를 구분할 수 있다.

\begin{dhkcode}
import asyncio
import inspect
from functools import wraps


def trace(func):

    if inspect.iscoroutinefunction(func):

        @wraps(func)
        async def async_wrapper(*args, **kwargs):
            print("enter async")
            return await func(*args, **kwargs)

        return async_wrapper

    @wraps(func)
    def sync_wrapper(*args, **kwargs):
        print("enter sync")
        return func(*args, **kwargs)

    return sync_wrapper


@trace
async def calculate(x):
    return x * 2


print(inspect.iscoroutinefunction(calculate))
print(asyncio.run(calculate(5)))
\end{dhkcode}

일반 Python 스크립트에서의 예상 출력은 다음과 같다.

\begin{dhkzsh}
True
enter async
10
\end{dhkzsh}

이 코드를 이미 이벤트 루프가 실행 중인 Jupyter 환경에서 실행한다면

\zsh{asyncio.run()} 대신 다음과 같이 호출해야 한다.

\begin{dhkcode}
result = await calculate(5)
print(result)
\end{dhkcode}

\begin{dhknote}
\begin{dhksmry}[centered]
제네레이터 함수
\end{dhksmry}
\end{dhknote}

제네레이터 함수도 비슷하게 실행 시점을 구별해야 한다.

\begin{dhkcode}
def trace(func):

    def wrapper(*args, **kwargs):
        print("wrapper called")
        return func(*args, **kwargs)

    return wrapper


@trace
def numbers():
    print("generator body")
    yield 10
    yield 20


iterator = numbers()

print(next(iterator))
\end{dhkcode}

예상 출력:

\begin{dhkzsh}
wrapper called
generator body
10
\end{dhkzsh}

왜 출력 순서가 이렇게 될까?

\zsh{numbers()}는 실제로 일반 wrapper 함수 객체를 호출한다.

따라서 \zsh{"wrapper called"}가 먼저 출력된다.

그다음 wrapper가 원본 제네레이터 함수를 호출하여 제네레이터 객체를 반환한다.

하지만 이 시점에는 원본 제너레이터 본문이 아직 실행되지 않았다.

\zsh{next(iterator)}가 호출될 때 비로소 제너레이터 본문이 실행을 시작한다.

따라서 함수 객체의 종류뿐 아니라 데코레이터가 반환한 객체의 종류와 실제 호출 결과도 구별해야 한다.

\subsubsection{실전 데코레이터 설계에서 발생하는 경계 조건}

지금가지의 내용을 바탕으로 실전에서 자주 발생하는 문제를 정리해 보자.

\begin{dhktable}
|문제 상황|	발생 원인|	설계 시 확인할 사항|
|:---|:---|:---|
|원본 함수의 반환값이 사라짐|	wrapper에서 return 누락|	return func(*args, **kwargs)가 필요한가?|
|원본 함수가 너무 일찍 실행됨	|func 대신 func()를 데코레이션 시점에 호출|	함수 객체와 호출 결과를 구별했는가?|
|여러 호출의 상태가 공유됨|	wrapper들이 같은 클로저 셀 참조|	상태를 호출별·인스턴스별·함수별 중 어디에 저장할 것인가?|
|예외가 사라짐	|except 또는 finally에서 예외 억제	|예외를 전파할 것인가, 변환할 것인가?|
|클래스 메서드 바인딩 실패|	데코레이터가 classmethod 객체를 직접 감쌈|	최종 클래스 속성의 타입은 무엇인가?|
|클래스 기반 데코레이터에서 self 누락|	반환 객체가 __get__을 제공하지 않음|	호출 가능성과 디스크립터 여부를 구별했는가?|
|FastAPI에서 매개변수 검증 실패	|wrapper 시그니처와 메타데이터 문제|	__wrapped__, 실제 인수 전달이 적절한가?|
|비동기 엔드포인트 실행 오류	|코루틴 함수를 동기 wrapper로 감쌈|	최종 함수가 코루틴 함수인가?|
|등록된 함수와 이름이 달라짐|	등록 이후 다른 wrapper 적용 또는 이름 재바인딩|	레지스트리는 어느 객체를 참조하는가?|
|호출 횟수 집계가 부정확함|	공유 상태 또는 동시 실행|	상태 변경의 원자성·동기화가 필요한가?|
\end{dhktable}

\begin{dhknote}
\begin{dhksmry}[centered]
공유 상태와 동시 실행
\end{dhksmry}
\end{dhknote}

앞의 \zsh{monitor}에서 \zsh{calls}는 클로저 셀을 통해 유지된다.

이 상태는 동일한 wrapper를 호출하는 모든 호출자가 공유한다.

예를 들어 \zsh{Base.preidct(...)}와 \zsh{Derived.predict(...)}는 같은 wrapper 함수 객체를 사용하므로 호출 횟수도 공유한다.

그런데 여러 스레드에서 동시에 호출하는 환경이라면 다음 연산의 동기화까지 고려해야 한다.

\begin{dhkcode}
calls += 1
\end{dhkcode}

이 표현식을 모든 Python 구현과 실행 환경에서 안전한 원자적 카운터 연산이라고 가정하면 안 된다.

정확한 동시 호출 횟수 집계가 필요하다면 잠김(Lock) 등 적절한 동기화 수단을 사용해야 한다.

또한 호출 횟수가 인스턴스별로 분리되어야 한다면 하나의 공유 클로저 셀에 저장하는 방식이 요구사항에 맞지 않을 수 있다.

결국 클로저는 상태를 유지할 수 있는 메커니즘을 제공하지만, 그 상태의 공유 범위와 동시성 정책은 프로그램 설계자가 결정해야 한다.

\subsubsection{데코레이터를 분석하는 최종 절차}

이제 처음 보는 데코레이터 코드도 다음 순서로 분석하면 된다.

\begin{dhksmry}
Python 데코레이터 실행 분석 절차
1. 정의가 실행되는 위치를 확인한다.
   모듈 본문인지, 클래스 본문인지, 함수 내부인지 확인한다. 
   그 위치에 따라 이름이 바인딩되는 네임스페이스가 달라진다.

2. 데코레이터 표현식의 평가 순서를 추적한다.
   데코레이터 팩토리가 있다면 언제 호출되고 무엇을 반환하는지 확인한다.

3. 원본 함수 객체의 생성과 데코레이터 적용을 구별한다.
   함수 본문이 실행된 것인지, 함수 객체만 생성된 것인지 구별해야 한다.

4. 각 데코레이터의 입력 객체와 반환 객체를 기록한다.
   함수 객체, wrapper, 디스크립터, 호출 가능한 인스턴스 중 무엇이 
   전달되고 반환되는지 확인한다.

5. 최종 이름 바인딩과 외부 등록 참조를 확인한다.
   함수 이름이 어떤 객체를 참조하는지, 레지스트리는 어떤 객체를 
   참조하는지 구별한다.

6. 속성 조회와 메서드 바인딩을 분석한다.
   클래스 또는 인스턴스를 통해 접근한다면 __get__, __self__, 
   __func__를 확인한다.

7. 실제 호출에서 인수와 클로저를 추적한다.
   args, kwargs, 자유 변수, 클로저 셀, nonlocal에 의한 
   상태 변화를 분석한다.

8. 반환·예외·최종 상태를 확인한다.
   반환값이 그대로 전달되는지, 예외가 전파되는지, 상태 변경이 남는지,
   비동기 실행이 필요한지 확인한다.
\end{dhksmry}

이 절차는 데코레이터뿐 아니라 콜백(Callback), 이벤트 핸들러(Event Handler), 의존성 주입(Dependency Injection), 클래스 기반 프레임워크의 실행 흐름을 분석할 때도 활용할 수 있다.

\subsubsection{전체 개념 통합}

지금까지 학습한 내용을 객체 모델 관점에서 다시 연결해 보자.

\begin{dhktable}
|단원|	핵심 개념|	최종적으로 설명할 수 있어야 하는 것|
|:---|:---|:---|
|A.1|	함수 객체와 이름 바인딩|	func와 func()의 차이|
|A.2|	중첩 함수와 렉시컬 스코프|	함수가 정의된 위치와 외부 이름 참조|
|A.3|	LEGB|	이름 탐색이 이루어지는 네임스페이스|
|A.4|	클로저와 자유 변수|	외부 함수 종료 후에도 참조가 유지되는 이유|
|A.5|	nonlocal, __closure__|	클로저 셀을 통한 상태 재바인딩|
|A.6|	데코레이터 변환식|	@decorator의 함수 객체 변환|
|A.7|	Wrapper와 인수 전달|	호출 중개, 반환값, 예외 처리|
|A.8|	functools.wraps|	메타데이터와 실제 함수 동작의 구별|
|A.9|	데코레이터 팩토리|	설정값을 클로저로 유지하는 구조|
|A.10|	다중 데코레이터|	평가·적용·호출 순서의 차이|
|A.11|	메서드 바인딩과 디스크립터|	클래스 속성 조회와 바운드 메서드 생성|
|A.12|	등록형 데코레이터|	함수 객체와 외부 레지스트리의 참조 관계|
|A.13|	실행 모델 통합|	전체 객체 생성·바인딩·탐색·호출·반환·상태 변화 추적|
\end{dhktable}

이번 학습 범위를 하나의 문장으로 정리하면 다음과 같다.

데코레이터는 함수 객체를 전달받아 어떤 처리를 수행하고b객체를 반환하는 Python의 일반적인 호출 구조이며, 그 결과는 이름 바인딩·클로저·디스크립터·외부 레지스트리·실제 호출 환경에 따라 서로 다른 실행동작으로 연결된다.

이 원리를 이해했다면 새로운 데코레이터를 만났을 때 단순히 문법을 암기할 필요가 없다.

어떤 객체가 어디에 저장되어 있고, 어떤 시점에 조회되고 호출되는지를 추적하면 동작을 설명할 수 있기 때문이다.
